% 地球系统数值模拟教育部重点实验室开放基金申请书。
% 坐标单位 bp，原点为纸张左上角，参数来自用户提供的四页官方 PDF。
% 固定表单用原生文字和矢量线条重建；报告正文使用可跨页的正常排版。
\RequirePackage{geometry}
\RequirePackage{fontspec}
\RequirePackage{xeCJK}
\RequirePackage{bensz-fonts}
\RequirePackage{tikz}
\RequirePackage{lastpage}
\RequirePackage{amsmath,amssymb}
\RequirePackage{graphicx}
\newif\ifCiesmBlank
\geometry{paperwidth=595.3bp,paperheight=841.9bp,
  left=90bp,right=90bp,top=80.69bp,bottom=80bp,headheight=0pt,headsep=0pt}
\setmainfont[Path={\BenszFontsDir},AutoFakeBold=2.5]{Kaiti.ttf}
\setCJKmainfont[Path={\BenszFontsDir},AutoFakeBold=2.5]{Kaiti.ttf}
\setCJKsansfont[Path={\BenszFontsDir},AutoFakeBold=2.5]{Kaiti.ttf}
\setCJKmonofont[Path={\BenszFontsDir}]{Kaiti.ttf}
\newCJKfontfamily\CiesmFooterSong[Path={\BenszFontsDir}]{SimSun.ttf}
\newfontfamily\CiesmFooterTimes[Path={\BenszFontsDir}]{TimesNewRoman.ttf}
\newfontfamily\CiesmBibliographyLatin[Path={\BenszFontsDir},Ligatures=TeX,AutoFakeSlant=.2]{TimesNewRoman.ttf}
\xeCJKsetup{PunctStyle=plain,CJKecglue=\hskip .25em}
\linespread{1}
\setlength{\parindent}{2em}
\setlength{\parskip}{0pt}
\setlength{\topskip}{12bp}
\pagestyle{empty}
\AtBeginDocument{\fontsize{12bp}{22bp}\selectfont}

\pgfkeys{/ciesm/.is family,/ciesm,
  special-name/.initial={}, title/.initial={}, applicant/.initial={},
  institution/.initial={}, unit-count/.initial={}, sex/.initial={},
  birth-date/.initial={}, professional-title/.initial={}, degree/.initial={},
  start-year/.initial={}, start-month/.initial={}, start-day/.initial={},
  end-year/.initial={}, end-month/.initial={}, end-day/.initial={},
  total-budget/.initial={}, expected-results/.initial={}, keywords/.initial={},
  issuer/.initial={}, issue-date/.initial={}}
\pgfkeys{/ciesm/.cd,
  faculty/.initial={}, address/.initial={}, phone/.initial={}, email/.initial={},
  mentor/.initial={}, english-name/.initial={}, talent/.initial={}, office/.initial={},
  homepage/.initial={}, education/.initial={}, employment/.initial={}}
\NewDocumentCommand{\CiesmSetup}{+m}{\pgfkeys{/ciesm/.cd,#1}}
\newcommand{\CiesmValue}[1]{\pgfkeysvalueof{/ciesm/#1}}
\newcommand{\CiesmDate}[1]{\CiesmValue{#1-year}年\CiesmValue{#1-month}月\CiesmValue{#1-day}日}
\newcommand{\CiesmUrlValue}[1]{\begingroup
  \edef\ciesmUrl{\CiesmValue{#1}}\expandafter\url\expandafter{\ciesmUrl}\endgroup}

% 与官方表单一致的页脚；总页数随正文自动增长。
\AddToHook{shipout/foreground}{%
  \begin{tikzpicture}[remember picture,overlay]
    \node[anchor=base,inner sep=0pt] at
      ([xshift=297.65bp,yshift=-789.77bp]current page.north west)
      {\CiesmFooterTimes\CiesmFooterSong\fontsize{9bp}{9bp}\selectfont
       第\hspace{4.5bp}\thepage\hspace{4.5bp}页\hspace{4.5bp}共\hspace{4.5bp}\pageref{LastPage}\hspace{4.5bp}页};
  \end{tikzpicture}%
}

\newenvironment{CiesmCanvas}{%
  \noindent\begin{tikzpicture}[remember picture,overlay]
  \begin{scope}[shift={(current page.north west)},x=1bp,y=-1bp,
    every node/.style={inner sep=0pt,outer sep=0pt}]
}{\end{scope}\end{tikzpicture}\par\null}
\newcommand{\CiesmText}[3]{\node[anchor=base west] at (#1,#2) {#3};}

% 固定信息表的填写区：按真实高度检查，过长即报错，避免静默越框。
\newsavebox{\CiesmFieldBox}
% TikZ 绘图区使用 nullfont，测量可编辑文本时须恢复普通排版环境。
\newcommand{\CiesmMeasure}[2]{%
  \setbox#1=\hbox{\begin{pgfinterruptpicture}#2\end{pgfinterruptpicture}}%
}
\newcommand{\CiesmMeasureCell}[2]{%
  \CiesmMeasure{\CiesmFieldBox}{\begin{minipage}{#1bp}
    \centering\normalfont\fontsize{10.45bp}{16bp}\selectfont #2\par
  \end{minipage}}%
}
% 长字段保持同一字号和行距，增加行高；留出 4bp 的上下内边距。
\newcommand{\CiesmCellGrowth}[4]{%
  \CiesmMeasureCell{#2}{#4}%
  \pgfmathsetmacro{#1}{max(0,ceil((\ht\CiesmFieldBox+\dp\CiesmFieldBox)/1bp+4-#3))}%
}
\newcommand{\CiesmCell}[5]{%
  \CiesmMeasureCell{#3}{#5}%
  \ifdim\dimexpr\ht\CiesmFieldBox+\dp\CiesmFieldBox\relax>#4bp
    \PackageError{ciesm-open-fund}{Information field exceeds its box}
      {Adjust this template's local field dimensions; do not truncate the value.}%
  \fi
  \node at ({#1+#3/2},{#2+#4/2}) {\usebox{\CiesmFieldBox}};
}
\newcommand{\CiesmCoverField}[5]{%
  \CiesmText{111}{#1}{#2}%
  \draw[line width=.75bp] (#3,{#1+3.2}) -- (#4,{#1+3.2});
  \CiesmMeasure{\CiesmFieldBox}{\normalfont\fontsize{15.95bp}{20bp}\selectfont\CiesmValue{#5}}%
  \ifdim\wd\CiesmFieldBox>\dimexpr#4bp-#3bp\relax
    \CiesmMeasure{\CiesmFieldBox}{\begin{minipage}{\dimexpr#4bp-#3bp\relax}
      \centering\normalfont\fontsize{12bp}{13bp}\selectfont\CiesmValue{#5}\par
    \end{minipage}}%
    \ifdim\dimexpr\ht\CiesmFieldBox+\dp\CiesmFieldBox\relax>27bp
      \PackageError{ciesm-open-fund}{Cover field '#5' exceeds two lines}
        {Adjust the local field layout; keep the complete value.}%
    \fi
    \node[anchor=south] at ({(#3+#4)/2},#1) {\usebox{\CiesmFieldBox}};
  \else
    \node[anchor=base] at ({(#3+#4)/2},#1) {\usebox{\CiesmFieldBox}};
  \fi
}

% 长课题名称保持封面字段字号，逐行画下划线；后续字段整体下移一行。
\newcommand{\CiesmCoverTitle}{%
  \def\ciesmCoverShift{0}%
  \CiesmMeasure{\CiesmFieldBox}{\normalfont\fontsize{15.95bp}{31.1bp}\selectfont\CiesmValue{title}}%
  \ifdim\wd\CiesmFieldBox>297bp
    \def\ciesmCoverShift{31.1}%
    \CiesmText{111}{337.49}{课题名称：}%
    \draw[line width=.75bp] (198,340.69) -- (495,340.69);
    \draw[line width=.75bp] (198,371.79) -- (495,371.79);
    \CiesmMeasure{\CiesmFieldBox}{\begin{minipage}[b]{297bp}
      \centering\normalfont\fontsize{15.95bp}{31.1bp}\selectfont\CiesmValue{title}\par
    \end{minipage}}%
    \ifdim\dimexpr\ht\CiesmFieldBox+\dp\CiesmFieldBox\relax>50bp
      \PackageError{ciesm-open-fund}{Cover field 'title' exceeds two lines}
        {Adjust the cover field width; do not shrink or truncate the title.}%
    \fi
    \node[anchor=base] at (346.5,368.59) {\usebox{\CiesmFieldBox}};
  \else
    \CiesmCoverField{337.49}{课题名称：}{198}{495}{title}%
  \fi
}

\newcommand{\CiesmMakeCover}{%
  \begin{CiesmCanvas}
    \fontsize{20.05bp}{31.2bp}\selectfont
    \node[anchor=base] at (297.65,185.69)
      {\bfseries\xeCJKsetup{CJKglue=\hskip .23bp}地球系统数值模拟教育部重点实验室开放基金};
    \node[anchor=base] at (297.65,216.89)
      {\bfseries\xeCJKsetup{CJKglue=\hskip .23bp}课题申请书};
    \fontsize{15.95bp}{31.1bp}\selectfont
    \CiesmCoverField{306.41}{专项名称：}{198}{495}{special-name}
    \CiesmCoverTitle
    \begin{scope}[shift={(0,\ciesmCoverShift)}]
    \CiesmCoverField{368.69}{课题申请人：}{208.5}{495}{applicant}
    \CiesmCoverField{399.77}{课题依托单位：}{225.75}{489.95}{institution}
    \CiesmText{111}{430.85}{起止年限：\CiesmValue{start-year}年\CiesmValue{start-month}月}
    \CiesmText{291}{430.85}{至}
    \CiesmText{323.04}{430.85}{\CiesmValue{end-year}年\CiesmValue{end-month}月}
    \fontsize{18bp}{23.04bp}\selectfont
    \node[anchor=base] at (297.65,533.21) {\CiesmValue{issuer}};
    \node[anchor=base] at (297.65,556.25) {\CiesmValue{issue-date}};
    \end{scope}
  \end{CiesmCanvas}
  \clearpage
}

\newcommand{\CiesmMakeInformation}{%
  \begin{CiesmCanvas}
    \CiesmCellGrowth{\ciesmTitleGrowth}{233}{19}{\CiesmValue{title}}%
    \CiesmCellGrowth{\ciesmResultsGrowth}{335}{39}{\raggedright\CiesmValue{expected-results}}%
    \CiesmCellGrowth{\ciesmKeywordsGrowth}{284}{20}{\CiesmValue{keywords}}%
    \pgfmathsetmacro{\ciesmTitleHeight}{19+\ciesmTitleGrowth}%
    \pgfmathsetmacro{\ciesmResultsHeight}{39+\ciesmResultsGrowth}%
    \pgfmathsetmacro{\ciesmKeywordsHeight}{20+\ciesmKeywordsGrowth}%
    \pgfmathsetmacro{\ciesmAbstractShift}{\ciesmTitleGrowth+\ciesmResultsGrowth}%
    \pgfmathsetmacro{\ciesmAbstractHeight}{288-\ciesmAbstractShift-\ciesmKeywordsGrowth}%
    \node[anchor=base,font=\fontsize{15bp}{20bp}\selectfont] at (297.65,93.05) {一、基本信息表};
    \draw[line width=1.44bp] (67.1,103.92) rectangle (528.2,610.12);
    \draw[line width=.72bp] (182.15,103.92) -- (182.15,{273.13+\ciesmAbstractShift});
    \foreach \x in {423.15,495.3}
      \draw[line width=.96bp] (\x,103.92) -- (\x,{124.91+\ciesmTitleGrowth});
    \fontsize{10.45bp}{16bp}\selectfont
    \CiesmText{103.56}{117.05+\ciesmTitleGrowth/2}{\bfseries 课题名称}
    \CiesmText{432.84}{117.05+\ciesmTitleGrowth/2}{\bfseries 参加单位数}
    \CiesmCell{186}{104.5}{233}{\ciesmTitleHeight}{\CiesmValue{title}}
    \CiesmCell{498}{104.5}{27}{\ciesmTitleHeight}{\CiesmValue{unit-count}}
    \begin{scope}[shift={(0,\ciesmTitleGrowth)}]
    \draw[line width=1.44bp] (67.1,145.87) -- (528.2,145.87);
    \foreach \y in {124.91,226.91}
      \draw[line width=.72bp] (67.1,\y) -- (528.2,\y);
    \foreach \y in {166.86,187.46}
      \draw[line width=.72bp] (97.85,\y) -- (528.2,\y);
    \draw[line width=.96bp] (67.1,208.18) -- (528.2,208.18);
    \draw[line width=.96bp] (67.1,{273.13+\ciesmResultsGrowth}) -- (528.2,{273.13+\ciesmResultsGrowth});
    \draw[line width=.72bp] (97.85,145.87) -- (97.85,208.18);
    \foreach \x in {289.05,423.15}
      \draw[line width=.72bp] (\x,124.91) -- (\x,208.18);
    \CiesmText{93}{137.69}{\bfseries 课题起始时间}
    \CiesmText{197.52}{137.69}{\CiesmDate{start}}
    \CiesmText{324.48}{137.69}{\bfseries 课题结束时间}
    \CiesmText{432.36}{137.69}{\CiesmDate{end}}
    \CiesmText{77.16}{163.61}{\bfseries 申}
    \CiesmText{77.16}{179.69}{\bfseries 请}
    \CiesmText{77.16}{195.65}{\bfseries 人}
    \CiesmText{119.04}{159.05}{姓\hspace{20.99bp}名}
    \CiesmText{345.6}{159.05}{性别}
    \CiesmText{119.04}{179.69}{出生日期}
    \CiesmText{345.6}{179.69}{职称}
    \CiesmText{119.04}{200.21}{最高学位}
    \CiesmText{335.04}{200.21}{所在单位}
    \CiesmText{99}{220.13}{\bfseries 申请经费情况}
    \CiesmText{187.56}{220.13}{课题总经费：}
    \draw[line width=.48bp] (250.55,221.25) -- (282,221.25);
    \CiesmText{282}{220.13}{万元}
    \CiesmText{103.56}{252.41+\ciesmResultsGrowth/2}{\bfseries 预期成果}
    \CiesmCell{186}{146.5}{99}{19}{\CiesmValue{applicant}}
    \CiesmCell{427}{146.5}{97}{19}{\CiesmValue{sex}}
    \CiesmCell{186}{167.2}{99}{19}{\CiesmValue{birth-date}}
    \CiesmCell{427}{167.2}{97}{19}{\CiesmValue{professional-title}}
    \CiesmCell{186}{187.9}{99}{19}{\CiesmValue{degree}}
    \CiesmCell{427}{187.9}{97}{19}{\CiesmValue{institution}}
    \CiesmCell{250.55}{208.5}{31.45}{17}{\CiesmValue{total-budget}}
    \CiesmCell{187}{230}{335}{\ciesmResultsHeight}{\raggedright\CiesmValue{expected-results}}
    \end{scope}
    \draw[line width=.96bp] (67.1,{588.43-\ciesmKeywordsGrowth}) -- (528.2,{588.43-\ciesmKeywordsGrowth});
    \draw[line width=.48bp] (234.55,{588.43-\ciesmKeywordsGrowth}) -- (234.55,610.12);
    \CiesmText{72.48}{284.21+\ciesmAbstractShift}{{\bfseries 课题摘要}（包括标点符号不超过400字）}
    \CiesmText{72.48}{599.45-\ciesmKeywordsGrowth/2}{\bfseries 关键词（用分号隔开，最多5个）}
    \pgfmathsetmacro{\ciesmKeywordsTop}{589-\ciesmKeywordsGrowth}%
    \CiesmCell{239}{\ciesmKeywordsTop}{284}{\ciesmKeywordsHeight}{\CiesmValue{keywords}}
    \CiesmMeasure{\CiesmFieldBox}{\begin{minipage}{444bp}
      \fontsize{12bp}{22bp}\selectfont\setlength{\parindent}{2em}
      \ifCiesmBlank\else\input{extraTex/abstract}\fi\par
    \end{minipage}}%
    \ifdim\dimexpr\ht\CiesmFieldBox+\dp\CiesmFieldBox\relax>\ciesmAbstractHeight bp
      \PackageError{ciesm-open-fund}{Abstract exceeds its box}
        {Adjust the information rows or shorten extraTex/abstract.tex; the limit is 400 characters.}%
    \fi
    \node[anchor=north west] at (76,{294+\ciesmAbstractShift}) {\usebox{\CiesmFieldBox}};
  \end{CiesmCanvas}
  \clearpage
}

% 正文不装入固定框，正常段落、公式、图表与参考文献可以跨页。
\newcommand{\CiesmReportTitle}[1]{%
  {\par\noindent\fontsize{14.05bp}{25.2bp}\selectfont #1\par}\nobreak
}
\NewDocumentCommand{\CiesmReportPart}{O{22bp} m m}{%
  {\par\noindent\fontsize{14.05bp}{#1}\selectfont #2{\fontsize{12bp}{22bp}\selectfont #3}\par}\nobreak
}
\newcommand{\CiesmPrompt}[1]{\par\noindent #1\par\nobreak}
\newcommand{\CiesmTemplateBreak}{\ifCiesmBlank\linebreak\fi}
\newcommand{\CiesmAnswer}[1]{\ifCiesmBlank\else\input{#1}\fi\par\vspace{22bp}}

% 填写稿采用真正的标题层级；空白入口保持官方四页的定位与留白。
% 官方两部分为 section，各部分的官方栏目为 subsection，栏目内为 1.1 等。
\newcommand{\CiesmApplicationMode}{%
  \RequirePackage{placeins}
  \RequirePackage{caption}
  \RequirePackage{ragged2e}
  % 参考文献著录与国自然口径统一：GB/T 7714 顺序编码制。
  % bensz-nsfc 走 BibTeX + gbt7714-nsfc.bst，本模板走 biber，故用 biblatex-gb7714-2015
  % 实现同一标准；sorting=none 保证按正文引用顺序编号，maxbibnames=3 对应“列三位后加等”。
  \RequirePackage[backend=biber,style=gb7714-2015,sorting=none,
    % maxcitenames 同时管 \fullcite；设为 3 才能保证代表作里申请人本人的署名可见。
    maxbibnames=3,minbibnames=3,maxcitenames=3,mincitenames=3,
    gbalign=gb7714-2015,gbpub=true,gbtype=true,doi=true,url=false,isbn=false]{biblatex}
  \ctexset{
    section={format=\normalfont\fontsize{14.05bp}{22bp}\selectfont,
      name={},number={},beforeskip=12bp,afterskip=6bp},
    subsection={format=\normalfont\fontsize{12bp}{22bp}\selectfont,
      name={},number={},beforeskip=10bp,afterskip=4bp},
    subsubsection={format=\normalfont\bfseries\fontsize{12bp}{22bp}\selectfont,
      number=\arabic{subsection}.\arabic{subsubsection},
      beforeskip=8bp,afterskip=2bp},
    paragraph={format=\normalfont\bfseries\fontsize{12bp}{22bp}\selectfont,
      number=\arabic{subsection}.\arabic{subsubsection}.\arabic{paragraph},
      runin=false,beforeskip=5bp,afterskip=0bp}}
  \setcounter{secnumdepth}{4}
  \setlength{\emergencystretch}{2em}
  \clubpenalty=10000
  \widowpenalty=10000
  \renewcommand{\thesection}{\arabic{section}}
  \renewcommand{\thesubsection}{\arabic{subsection}}
  \RenewDocumentCommand{\CiesmReportPart}{O{22bp} m m}{%
    \FloatBarrier\section{##2{\fontsize{12bp}{22bp}\selectfont ##3}}}
  \renewcommand{\CiesmPrompt}[1]{\FloatBarrier\subsection{##1}}
  \renewcommand{\CiesmAnswer}[1]{\input{##1}\par}
  \DeclareCaptionFont{ciesm}{\fontsize{10.5bp}{16bp}\selectfont}
  \captionsetup[figure]{font=ciesm,labelfont=normalfont,labelsep=quad,
    skip=6bp,singlelinecheck=false,justification=justified}
  \setlength{\textfloatsep}{12bp}
  \setlength{\intextsep}{12bp}
  % 参考文献使用 10.5bp 字号、16bp 行距，条目之间不额外留白。
  \setlength{\bibitemsep}{0pt}
  \setlength{\bibparsep}{0pt}
  % 楷体不含全部拉丁附加符号；题录西文使用共享 Times 字体保全作者姓名。
  \renewcommand{\bibfont}{\normalfont\CiesmBibliographyLatin\fontsize{10.5bp}{16bp}\selectfont}
  % 编号贴正文左边界，续行保留悬挂缩进；左对齐避免长英文词拉大字距。
  \AtBeginBibliography{\RaggedRight\setlength{\parindent}{0pt}\setlength{\parskip}{0pt}}
  % 文章号不再在此处理：gb7714-2015 的 sourcemap 已把 article 的 eid 映射进 pages
  % （且不覆盖原有 pages），旧的 numeric 样式下用来去重的 \clearfield{pages} 在这里
  % 会反过来把文章号整个抹掉，故移除。
  \defbibheading{ciesm}{\subsubsection*{主要参考文献}}
}
\newcommand{\CiesmFullCite}[1]{{\CiesmBibliographyLatin\fullcite{#1}}}

% 图和图注作为一个浮动体，限宽限高且始终保留纵横比，不裁剪原 PDF。
\newcommand{\CiesmFigure}[3]{%
  \begin{figure}[!htbp]
    \centering
    \includegraphics[width=\linewidth,height=.62\textheight,keepaspectratio]{#1}
    \caption{#2}\label{#3}
  \end{figure}
}

% 预算表保留官方列宽与底色；概要说明换行时相应行自动增高。
\definecolor{CiesmBudgetGray}{gray}{.953}
\newsavebox{\CiesmBudgetNameBox}
\newsavebox{\CiesmBudgetAmountBox}
\newsavebox{\CiesmBudgetNoteBox}
\newdimen\CiesmBudgetHeight
\newcommand{\CiesmBudgetRow}[4]{%
  \CiesmMeasure{\CiesmBudgetNameBox}{\parbox{170bp}{\normalfont\fontsize{9.95bp}{14bp}\selectfont\raggedright #2\par}}%
  % 空白入口复用同一科目表，但不展开填写稿的金额与说明。
  \CiesmMeasure{\CiesmBudgetAmountBox}{\parbox{55bp}{\normalfont\fontsize{9.95bp}{14bp}\selectfont\centering\ifCiesmBlank\else #3\fi\par}}%
  \CiesmMeasure{\CiesmBudgetNoteBox}{\parbox{93bp}{\normalfont\fontsize{9.95bp}{14bp}\selectfont\centering\ifCiesmBlank\else #4\fi\par}}%
  \CiesmBudgetHeight=20.6bp
  \foreach \boxname in {\CiesmBudgetNameBox,\CiesmBudgetAmountBox,\CiesmBudgetNoteBox}{%
    \ifdim\dimexpr\ht\boxname+\dp\boxname+6bp\relax>\CiesmBudgetHeight
      \global\CiesmBudgetHeight=\dimexpr\ht\boxname+\dp\boxname+6bp\relax
    \fi
  }%
  \pgfmathsetmacro{\ciesmRowBottom}{\ciesmRowTop+\CiesmBudgetHeight/1bp}
  \fill[CiesmBudgetGray] (104.15,\ciesmRowTop) rectangle (323.9,\ciesmRowBottom);
  \draw[line width=.72bp] (104.15,\ciesmRowBottom) -- (491.15,\ciesmRowBottom);
  \node at (120.825,{(\ciesmRowTop+\ciesmRowBottom)/2}) {#1};
  \node[anchor=west] at (149.04,{(\ciesmRowTop+\ciesmRowBottom)/2}) {\usebox{\CiesmBudgetNameBox}};
  \node at (355.8,{(\ciesmRowTop+\ciesmRowBottom)/2}) {\usebox{\CiesmBudgetAmountBox}};
  \node at (439.425,{(\ciesmRowTop+\ciesmRowBottom)/2}) {\usebox{\CiesmBudgetNoteBox}};
  \let\ciesmRowTop\ciesmRowBottom
}
\newenvironment{CiesmBudget}{%
  \begin{CiesmCanvas}
    \CiesmText{90}{86.69}{\fontsize{14.05bp}{22bp}\selectfont 三、课题经费预算表}
    \fontsize{9.95bp}{14bp}\selectfont
    \fill[CiesmBudgetGray] (104.15,110.56) rectangle (491.15,151.76);
    \draw[line width=.72bp] (104.15,110.56) -- (491.15,110.56);
    \draw[line width=.72bp] (137.5,131.16) -- (491.15,131.16);
    \draw[line width=.72bp] (104.15,151.76) -- (491.15,151.76);
    \CiesmText{110.76}{134.69}{\bfseries 序号}
    \CiesmText{200.52}{124.49}{\bfseries 预算科目名称}
    \CiesmText{325.68}{124.49}{\bfseries 经费（万元）}
    \CiesmText{419.4}{124.49}{\bfseries 概要说明}
    \CiesmText{218.16}{145.01}{（1）}
    \CiesmText{343.2}{145.01}{（2）}
    \CiesmText{426.84}{145.01}{（3）}
    \def\ciesmRowTop{151.76}
}{%
    \draw[line width=.72bp] (137.5,110.56) -- (137.5,\ciesmRowTop);
    \pgfmathsetmacro{\ciesmRowBottom}{\ciesmRowTop+20.6}
    \fill[CiesmBudgetGray] (104.15,\ciesmRowTop) rectangle (323.9,\ciesmRowBottom);
    \CiesmText{115.68}{\ciesmRowTop+13.85}{共计}
    \node at (355.8,{\ciesmRowTop+10.3}) {\CiesmValue{total-budget}};
    \CiesmText{436.92}{\ciesmRowTop+13.85}{/}
    \foreach \x in {104.15,323.9,387.7,491.15}
      \draw[line width=.72bp] (\x,110.56) -- (\x,\ciesmRowBottom);
    \draw[line width=.72bp] (104.15,\ciesmRowBottom) -- (491.15,\ciesmRowBottom);
    \ifdim\ciesmRowBottom bp>740bp
      \PackageError{ciesm-open-fund}{Budget table exceeds one page}
        {Shorten the budget notes in extraTex/budget.tex.}%
    \fi
  \end{CiesmCanvas}
}
